\pdfminorversion=7
\documentclass[draft,linenumbers]{agujournal2025}
\usepackage{booktabs,amsmath,amssymb,microtype,url,graphicx,array}
\setlength{\emergencystretch}{3em}
\input{review_completion_claims.tex}
\input{simple_research_claims.tex}
\begin{document}
\journalname{Water Resources Research}
\title{Incomplete river-temperature years: separating recovery, screening and trend inference}
\authors{[Author names]\affil{1}}
\affiliation{1}{[Affiliations]}
\authoraddr{[Corresponding author and email]}
\authorrunninghead{Anonymous}\titlerunninghead{Incomplete river-temperature years}
\keypoints
{Observed windows can disprove a complete-year-only heat record, but most witnesses pass a 90 percent warm-season rule}
{Distribution regression and calibrated stochastic reconstruction have similar annual accuracy in difficult gaps}
{Retaining years improves the evaluated finite-record trends; forward interval coverage and useful release remain limited}
\maketitle\par
\begin{abstract}
Incomplete river-temperature years can contain information about sustained heat, but the evidence needed to recover an annual value, release it, or infer a trend differs. At \RRecordStations{} of \RRecordDenominator{} long-record stations, fully observed 30-day windows in \RRecordStationYears{} incomplete years exceed every complete-year heat maximum. Of these witnesses, \WAWarmAdmitted{} pass a station-specific 90\% warm-season rule, so the observation-only result mainly challenges strict annual completeness. Across twenty source repetitions, annual distribution regression and equally calibrated stochastic reconstruction have similar errors on difficult gaps. We jointly degrade multiyear archives at \WAStations{} stations before local refitting and compare retention and estimator effects on the same finite-record slopes. Bounded annual distributions use a stated tail extrapolation and identical marginals across three joint-rank constructions. For controlled gaps, nominal 95\% D2 trend intervals cover \WAOfflineCoverage\% offline and \WAForwardCoverage\% in retrospective forward evaluation. Task-level calibration is available for \WATaskForwardPercent\% of forward controlled-gap cases under empirical support checks, with no population coverage guarantee. A three-output rule retains bounds or withholds values, but its individual classifier score does not certify group risk. Exploratory Canadian comparisons use published daily means or reported-hour records with unresolved quality or time basis; independent strict-reference validation remains open. These findings support preserving observed heat information while restricting claims about numerical recovery, transfer and trend uncertainty.
\begin{plainlanguagesummary}
Gaps in a river's water-temperature record do not necessarily erase information about a hot year. If thirty consecutive observed days are already warmer than the record among fully observed years, the complete-year-only record is incomplete. Most such examples in our archive would also be retained by a less restrictive warm-season rule. We compared several ways to estimate missing values and then tested what they imply for a trend across many years. A calibrated classical reconstruction was about as accurate as a more elaborate annual model. Trend intervals were less reliable when applied to later years, and screening rules did not provide a general safety guarantee. The analysis distinguishes observed facts, estimated values and cases where the evidence is insufficient.
\end{plainlanguagesummary}
\end{abstract}

\section{Introduction}
Stream-temperature records support studies of thermal habitat, climate change and water management \cite{caissie2006,michel2020,vanhamel2024}. The annual maximum of a 30-day mean summarizes sustained heat and is related to the thermal peak used in river-network studies \cite{beaufort2022}. A monitoring gap can affect this statistic very differently depending on its timing, the observations remaining in the year and the information available from other years or stations.

Prior work has already shown that partial stream-temperature records can support some annual analyses. \citeA{johnson2021} evaluated missing intervals and imputation for annual mean, phase and amplitude, finding limitations that varied among rivers. \citeA{beaufort2022} used summer observations and climate correction for a multiyear thermal-peak metric. Neither should be interpreted as requiring every calendar year to be complete. Our question is narrower: what do incomplete daily records support about an annual sustained-heat maximum, numerical publication, and a finite-record trend? We compare strict completeness with a station-specific warm-season rule and distinguish a mathematical observation bound from model-based estimation.

Regression record augmentation and maintenance-of-variance methods have a long hydrological history \cite{matalas1964,hirsch1982,vogel1985}. A point prediction, a predictive distribution and a decision to release a value have different targets. Distributional imputation should be assessed for the downstream quantity rather than only for daily squared error \cite{naf2023,gneiting2007}. Here we compare a daily plug-in, simple corrections, calibrated stochastic reconstruction and annual distribution regression. We then jointly hide portions of a station's multiyear record before local refitting and evaluate whether retaining years changes the slope of that finite record. Finally, we examine whether the annual distributions support trend uncertainty and whether visible information can identify cases that should not receive a numerical estimate.

The hypotheses are deliberately conditional: target-oriented recovery can improve on a daily plug-in without necessarily beating a calibrated stochastic comparator; retaining years can improve a finite-record slope under a specified degradation mechanism; and annual and trend uncertainty require separate evaluation. The evidence map in Table~\ref{tab:evidence-map} identifies the comparisons and limits.
\input{claim_evidence_table.tex}

\section{Data and study design}
The source archive contains previously assembled European and US daily observations: 581 physical stations, 185 connected-record groups, 27 regions and 6,170 unique complete station-years. NASA POWER daily air temperature is the common meteorological input \cite{nasapower2026}. Artificial controlled gaps of 30, 60 or 90 days are crossed with overlap with a training-period climatological peak window; empirical incomplete calendars provide a separate family. Each of twenty source repetitions redraws masks and group-role composition. In the leave-one-year-out source comparison, fitting, interval calibration, score learning and rule calibration use disjoint connected groups outside each held region. These repetitions describe variability within the available archive, not twenty independent climates. Figure~\ref{fig:architecture} separates the observation-only, single-year and whole-archive evidence.

We use complete daily station-years as reference outcomes for artificial masking. Water observations hidden by a mask never enter the local fit for that case. Features available when a value is evaluated include the visible gap geometry, a training-period climatological peak, local water history, air conditions and available neighboring stations. The true peak location and hidden water values are evaluation quantities. Observation quality is inherited from the assembled source records; a reported value is not assumed to be error-free physical truth. An observed lower bound is valid for the reported series if its thirty measurements are comparable and valid.

The annual target is
\begin{equation}
 H30_y=\max_{j}\frac{1}{30}\sum_{d=j}^{j+29}Y_{yd}.
\end{equation}
For an incomplete year, the largest fully observed 30-day mean $L_y$ satisfies $H30_y\ge L_y$ for the same reported series. For 1981--2025 we select stations with at least fifteen complete reference years and compare every incomplete-year bound with the largest complete-year H30. We also evaluate whether those same witnesses pass a 95\% annual-coverage rule and a 90\% station-specific warm-season rule. The five consecutive warm months are obtained from the station's climatology after excluding the evaluated year. These are explicit study comparators; the seven-week missing-period result of \citeA{johnson2021} concerns annual signal parameters, not a validated H30 record rule.
\input{architecture_figures.tex}

\section{Methods}
\subsection{Annual recovery and decision to release}
The daily plug-in computes H30 after a selected daily reconstruction. Functional debiasing corrects its annual value, whereas D2 learns annual-error quantiles using visible features and simulated trajectories. A stochastic regression comparator generates conditional residual bridges and calculates H30 for each realization. D2 and the stochastic comparator have the same source interval-calibration roles and bootstrap multiplicities. Point estimates are medians; MAE, CRPS, 90\% interval coverage and interval score are reported on available cases, with paired comparisons on common cases.

In the twenty historical source fits, we apply the same observed-window support projection to both methods after their original calibration. The seven stored quantiles do not determine their outer tails; for comparable CRPS integration we extend the tails using adjacent normal-quantile spacings and enforce $L_y$ when available. The original model weights and calibration offsets remain fixed in this source re-analysis. In the whole-archive study described below, offsets for both methods are instead re-estimated on their original independent calibration roles after projection. Neither approach makes the assumed tail shape a verified full distribution.

The conditional-mean mechanism is useful but limited. For observed information $X$ and 30-day averaging vectors $w_j$, convexity gives
\[
 \max_j w_j^\top E(Y\mid X)\le E\{\max_j w_j^\top Y\mid X\}.
\]
It concerns the true conditional mean and conditional expectation; an approximate learner need not underestimate every individual year. Calibration, local adaptation, model bias and the chosen point summary remain empirical questions.

We evaluate a rule that reports an estimated value, an observed lower bound or an unknown state. Numerical release requires an available D2, at least four complete historical years, visible features within source-derived limits and a raw individual classifier score no greater than 0.10. The score was trained for the original unprojected D2 and is not a calibrated group-risk certificate for the projected value. A seven-day close match between visible water and air temperature is tested for the imposed exact-air-substitution scenario; that detector is not evaluated for other stresses. No observed bound is released when this quality flag is raised. The primary factorial trend comparison is separate. A second exploratory analysis passes released values and bounds into a Gaussian working likelihood, treating released numerical estimates as exact and each bound $L$ as a right-censored response with contribution $P(H30\ge L)$. Unknown cases remain in availability denominators. The resulting regression interval includes working-model variability and is not a task-calibrated missingness-only interval.

\input{simple_research_methods.tex}

Two small stress studies select one long-record station per region and five mask repetitions. One applies the target's exact masks simultaneously to all observed donor stations; the other changes annual missingness probabilities to 0.90 or 0.18 according to whether the year's air H30 exceeds the first-seven-year 80th percentile. The latter changes the realized masks and missing-year count as well as the heat association. A station with insufficient air data is unavailable; neither study constitutes a new climatic replication.

\section{Results}
\subsection{What the observations establish}
\input{figure_observed_records.tex}
At \RRecordStations{} of \RRecordDenominator{} eligible stations, \RRecordStationYears{} incomplete years contain an observed window above the maximum complete-year H30 (Figure~\ref{fig:observed-records}). Six stations have a witness exceeding the complete-year record by more than 0.5$^{\circ}$C. Another incomplete year could still be warmer, and measurement error or a changed sensor location can affect physical interpretation.

The practical consequence depends on the rule. Of the 35 witnesses, \WAWarmAdmitted{} pass the station-specific 90\% warm-season criterion and 31 pass 95\% annual coverage. Among those rejected by the warm-season criterion, \WAWarmExcludedExceed{} exceed the largest \emph{observed window} among rule-accepted years. A rule-accepted incomplete year may itself have an unobserved higher H30, so that comparison concerns reported windows rather than a proven full-record ordering. The evidence therefore establishes a cost of strict complete-year selection; it does not attribute all 35 omissions to the more permissive rule.

\input{simple_external_results.tex}

\subsection{Annual point and distribution comparisons}
On controlled gaps overlapping at least half the training climatological peak, the median source-repetition H30 MAE is approximately 0.428$^{\circ}$C for the daily plug-in, 0.355 for functional debiasing, \WADTwoMAE{} for D2 and \WAStochasticMAE{} for calibrated stochastic regression (Figure~\ref{fig:estimators}). Median 90\% interval coverage is \WADTwoCoverage\% and \WAStochasticCoverage\%, respectively, with CRPS \WADTwoCRPS{} and \WAStochasticCRPS{}. On common cases D2 improves CRPS in \WADTwoCRPSWins{} of twenty repetitions and upper-tail CRPS in \WADTwoTailWins{}; equality of the methods is not established. Both distributions use the same support projection and tail approximation in this comparison. The practical evidence is a gain over the simple plug-in, not a demonstrated universal need for D2.

\input{uncertainty_paired_values.tex}
\input{support_projection_values.tex}
\input{peak_visibility_values.tex}
\input{annual_interval_expansion_values.tex}
\input{nova_scotia_exploratory_values.tex}
\input{nova_clock_sensitivity_values.tex}
\input{region_holdout_values.tex}
\input{au_prior_exposure_values.tex}
\input{source_rights_values.tex}
\input{open_source_values.tex}
\input{open_source_annual_values.tex}
\input{single_year_calibration_values.tex}
\input{single_year_trend_values.tex}
\input{strict_region_values.tex}
\input{input_precision_values.tex}
\input{release_matched_values.tex}
\input{release_regional_values.tex}
\input{bc_exploratory_values.tex}
\input{eccc_red_values.tex}
\input{simple_research_results.tex}
\input{simple_research_tables.tex}

A seed-zero audit of \RHModels{} source model bundles found no held test station or held test connected group in any source role. However, the forward HUC17 estimator used \RHHucCases{} cases from \RHHucStations{} other training-only station in \RHHucGroups{} connected group within the same named region, and the forward NVE interval-calibration role used \RHNveCases{} cases from \RHNveGroups{} other same-region groups. The former affects interpretation of \RHArchiveCases{} forward whole-archive station/repetition settings at \RHArchiveStations{} HUC17 stations; the two affected regions account for \RHFutureCases{} rows in the saved future-year release output. These are not held-target observations, but the forward source models are not uniformly trained outside the named target region.

A post-hoc strict-region sensitivity removes every source group appearing in each affected region, retains the other original role assignments and refits the two affected models. For forward controlled gaps, the common whole-archive comparison includes \SRPairN{} station/repetition cases at \SRPairStations{} stations; original and strict D2 trend RMSE are \SROldDTwo{} and \SRNewDTwo{}, and stochastic RMSE are \SROldStoch{} and \SRNewStoch{} $^{\circ}$C per decade. Across \SRFutureN{} forward controlled-gap annual cases, the three-output publication rate changes from \SROldRet\% to \SRNewRet\%, while failure among published values changes from \SROldRisk\% to \SRNewRisk\%. These selected sets differ. On \SRIntN{} common forward controlled-gap archives, release followed by the working likelihood yields RMSE \SRIntOld{} and \SRIntNew{}$^{\circ}$C per decade, with average working interval widths \SRIntOldWidth{} and \SRIntNewWidth{}. Among \SRAnnualN{} annual controlled-gap cases where the complete reference exceeds the visible bound, D2 nominal 90\% coverage changes from \SRAnnualOld\% to \SRAnnualNew\%. These are outcome-defined, repeated-mask comparisons. This smaller-source-sample sensitivity does not replace a strict-region rerun of every reported analysis or establish independent-climate transfer.

A further post-hoc source-rights sensitivity refits all 13 source models using seven candidate providers with recorded reuse paths, reassigns the source roles, and excludes all groups occurring in the target region. It reuses the original target masks and daily features. Among \OpenSourcePairs{} common forward controlled-gap station--repetition cases in \OpenSourceGroups{} connected groups, D2 and calibrated stochastic trend RMSE are \OpenSourceOldDTwo{} and \OpenSourceOldClassic{} $^{\circ}$C per decade under the original source models, versus \OpenSourceNewDTwo{} and \OpenSourceNewClassic{} under the candidate models. On a different four-cell factorial cohort of \OpenSourceFactorCases{} cases, the symmetric retention share for D2 is \OpenSourceOldRetention\% originally and \OpenSourceNewRetention\% for the candidate; these are ratios of mean contributions, not universal causal effects. On \OpenSourceCoverageCases{} common interval cases, nominal 95\% D2 coverage is \OpenSourceOldCoverage\% and \OpenSourceNewCoverage\%, respectively. Source providers, role assignments and strict region exclusion change together, so these shifts cannot be attributed to one factor. The analysis is retrospective and does not establish independent-climate transfer or permission to release the derived models and case data.

In the small synchronized-donor study, D2's forward controlled-gap RMSE changes from \WADonorBase{} to \WADonorStress{}$^{\circ}$C per decade on 30 matched station/repetition cases with identical target masks. In the heat-associated study, the offline controlled-gap comparison changes from \WAHeatBase{} to \WAHeatStress{} on 30 available paired station/repetition cases; five cases from one station are unavailable because its air history is insufficient. Different year masks and missing rates prevent attributing the latter difference solely to the heat association. Both studies have few stations and are diagnostic boundaries rather than population estimates. Under strict-region source sensitivity, the same controlled-gap donor comparison changes from \SRScenarioBaseNew{} to \SRDonorNew{}$^{\circ}$C per decade on \SRScenarioN{} common cases; the corresponding heat-associated stress RMSE is \SRHeatNew{} on that common cohort.

\subsection{External qualification and stress boundaries}
A subsequent Australian evaluation of a different model pipeline used the same public source archive \cite{bom2026}. Of \AURefYears{} station-years previously identified as candidate references, \AUFirstEvaluable{} entered that evaluation, including all \AUFormerTestEvaluable{} years previously designated for a separate F2 test. The historical F2 intake recorded its own formal test as closed, and this audit did not run it; the records are nevertheless not unopened targets for a new independent evaluation of the present pipeline. The earlier sample-support criterion allowed 20 observed hours per day, so it does not establish a strict complete 24-hour reference.

Nine Greater Wellington sites yielded no complete reference hot years under the original quality-code-600 criterion. The provider describes QC200 as unverified \cite{greaterwellington2016quality}, while the meanings of other returned legacy codes remain unresolved. A separate public NEON PRT archive \cite{marzolf2025neon} supplied half-hourly records at 24 river or stream sites. Using one fixed location code per site, zero sites had four complete years under either all 48 quality-passing half-hours per UTC day or a 46-of-48 sensitivity. Location codes do not establish identical hardware over time. The latter daily sensitivity is an observed-interval average and is not an exact complete day. An additional USGS intake of eight metadata-selected sites also produced no eligible four-year reference after prior-contact screening and daily quality checks. Five metadata-selected R\'eunion sites supplied Hub'Eau records \cite{hubeautemp2026} with quality code 1, but none supplied four complete July--June years after requiring 24 distinct valid hours per reported day; the provider-day time-zone convention remains unresolved. These qualification outcomes are not evidence that the temperatures are wrong, and they yield no external reconstruction-accuracy estimate.

We acquired \ECCCFiles{} ECCC Red River at Emerson files \cite{ecccautomated2024}; the site has a USGS alias \cite{usgsred2025} and shares a river with an upstream source station, so it is not an independent basin. The provider dictionary describes validation status V and quality grades 1--4, whereas \ECCCValidatedMissingGrade{} rows use the word VALIDATED without a grade and \ECCCUnresolvedIdentity{} rows use a station identifier different from the catalog. Under the fixed literal V and numeric-grade-1/2 rule, \ECCCCompleteYears{} complete local-clock years qualified; the daylight-saving convention remains unresolved. These exclusions do not establish that the other observations are erroneous, and no recovery accuracy was scored.

We also examined all 23 published British Columbia station records \cite{moore2024bc,callahan2025air2stream}; \BCStations{} stations had at least four calendar years with a finite published value every day. The Water Survey of Canada describes water temperature as an auxiliary output without standardized quality assurance \cite{wsc_temperature_faq}; the source study subsequently screened hourly measurements using spot checks and water--air comparisons and accepted daily means based on at least 20 hourly observations. The released derivative does not report the count and quality of each individual hour, so a complete year of published daily values is not a strict 24-hour-per-day reference. On \BCYears{} station-years with leave-one-year-out local fitting and \BCCases{} controlled-gap masks, the plug-in, D2 and calibrated stochastic H30 MAEs are \BCPluginMAE{}, \BCDTwoMAE{} and \BCStochasticMAE{}$^{\circ}$C, respectively; D2 and stochastic nominal 90\% coverage are \BCDTwoCoverage\% and \BCStochasticCoverage\%. These case-weighted summaries are exploratory: masks repeat within years, three sites lie in the Fraser River basin (Nechako, Thompson and Fraser), the source investigators selected their 23 sites for data availability, the published daily time-zone convention remains unresolved, and NASA air forcing uses coordinates rounded to 0.01 degree. The reference QC used water--air comparisons, so it does not independently test against that failure mechanism. This is a new station-record comparison against a weaker published-daily reference, not strict independent-climate or natural-outage validation.

A separate Nova Scotia exploratory amendment selected two stations and eight station-years after examining reported-clock completeness and agreement between two provincial releases \cite{novascotiawater2026}, but before calculating model prediction errors. Each reference day averages 24 finite, uniquely labelled hourly values; individual quality and estimated-value flags, the water time basis, and its alignment with local-solar-time NASA forcing remain unresolved. On \NSExpCases{} paired controlled-gap masks, D2 and calibrated stochastic H30 MAE are \NSExpDTwoMAE{} and \NSExpClassicMAE{}$^{\circ}$C; nominal 90\% coverage is \NSExpDTwoCover\% and \NSExpClassicCover\%, with widths \NSExpDTwoWidth{} and \NSExpClassicWidth{}$^{\circ}$C. These are repeated masks at \NSExpStations{} stations and \NSExpYears{} station-years, with single-year leave-out and three other complete years available for local fitting. A post-hoc sensitivity refits the local stages under reported-day boundaries shifted by $-1$ or $+1$ hour while holding masks, source model, calibration and forcing fixed. The maximum change in the eight reported reference H30 values is \NSClockHMax{}$^{\circ}$C, while the largest D2 point-prediction change over the \NSClockCases{} paired cases is \NSClockDTwoMax{}$^{\circ}$C; this does not establish the actual time basis or the cause of model sensitivity. This after-the-fact availability and source-agreement selection is not strict reference qualification, independent-climate confirmation, whole-archive degradation or a natural-outage test.

A further metadata-fixed USGS preflight excluded the known source, prior panels, downloaded caches and an API pilot station whose annual means were inspected and which had also been scored in an earlier E10 analysis. Of five remaining candidates with prior network screening, none had four calendar years of complete approved daily counts in 2010--2024 according to the USGS annual statistics service \cite{usgsstatistics2026}. The count screen may include estimated values and does not establish subdaily quality; no daily reconstruction was scored for these five stations. This availability-selected exploratory search does not describe the proportion of qualifying USGS stations overall.

An initial 15-station USGS intake illustrates why the access inventory matters: all 15 had already been processed in the national panel and ten had been scored in an earlier external comparison. They are opened historical evidence, not new independent validation. The subsequent eight-station roster was selected after excluding those contacts and other known source, evaluation and cached-value records. The prior-access and river-network screens are best effort; neither they nor different HUC identifiers prove independent weather or errors. All these intakes are exploratory and unregistered.

A complementary clean retrospective comparison selects the score-only cutoff in the held-region model's disjoint source rule-calibration groups, approximately matching the three-output numerical rate in those groups; it uses neither the target outcomes nor the target release count. On held-region empirical calendars, the three-output and source-calibrated score-only rules publish \SRCThreeRetention\% and \SRCScoreRetention\%, with failure rates \SRCThreeRisk\% and \SRCScoreRisk\%, respectively. On controlled gaps their publication rates are \SRCGapThreeRetention\% and \SRCGapScoreRetention\%, with failure rates \SRCGapThreeRisk\% and \SRCGapScoreRisk\%. The target retention rates differ, and this one-seed experiment does not establish a population-level tradeoff or group safety.

In the controlled-gap comparison, \RGGapPositive{} of \RGGapDefined{} regions with numerical outputs from both rules have a higher conditional error rate for three outputs, while \RGGapNegative{} have a lower rate; \RGGapUndefined{} region has undefined error because three outputs publishes none. These regional signs do not compare equal retention. For example, at the Finnish sites the three-output and score-only rates are \RGSykeThreeRisk\% and \RGSykeScoreRisk\%, while numerical retention is \RGSykeThreeRetention\% and \RGSykeScoreRetention\%. On the \RGCommonGroups{} of \RGAllGroups{} connected groups where both rules publish, the equal-group mean calendar-case error rates are \RGThreeGroupRisk\% and \RGScoreGroupRisk\%. Dropping each evaluation region in turn leaves the pooled error difference positive, with all models and source thresholds frozen; this is not a refitted leave-region-out study or a set of independent significance tests.

A retrospective comparison holds the numerical-output count fixed on each evaluated feature distribution. In clean empirical-calendar cases, both rules publish \MRCount{} values; three outputs have \MRThreeBad{} failures (\MRThreeRisk\%), while ranking cases only by the uncalibrated individual score yields \MRScoreBad{} (\MRScoreRisk\%). For forward controlled gaps, the corresponding \MRFutureCount{} values have failure rates \MRFutureThree\% and \MRFutureScore\%. Score ties use case identifiers, not outcomes. This comparator chooses its threshold from the evaluated score distribution and the three-output count, so it is a transductive descriptive benchmark rather than a prospectively calibrated policy. Equal numbers of numerical outputs do not equalize information retained as bounds or unknown states. In the imposed air-exposure calendar, only \MRAirCount{} numerical values remain; the two methods have \MRAirThreeBad{} and \MRAirScoreBad{} failures, respectively. These case-dependent results offer no group-risk guarantee.

The release comparison also exposes limits. In clean empirical-calendar cases, the fixed-score rule publishes about 94.7\% of values at 0.7\% observed error, while the three-output rule publishes about 85.5\% at a similar error rate. Under the imposed air-exposure stress, it publishes almost no numerical values. Hot-year and later-year stresses remain challenging despite source-range checks. These synthetic scenarios have no inferred field prevalence, and withholding almost every result is not a useful recovery guarantee.

\section{Discussion}
For the observation-only result, \IPDays{} finite daily values in \IPPanels{} selected project input panels agree with their float32 archive casts; the maximum rounding difference is \IPRounding{}$^{\circ}$C. No witness identity or 0.5$^{\circ}$C margin classification changes when the selected full-precision panels replace the canonical archive. These panels are themselves processed inputs and do not establish sensor accuracy.

The strongest observation result concerns the choice of reference archive. A fully observed window can establish that a strictly complete-year-only record misses a reported heat event without requiring an imputer. However, most witnesses would survive a less restrictive warm-season screen. The comparison with existing annual-signal and thermal-peak studies therefore matters: the result does not show that every partial-year practice is unnecessarily restrictive. It motivates reporting exactly which data and record statistic a rule excludes.

The whole-archive design shows that preserving supported years can improve the slope of the particular complete record used as reference. The estimate is conditional on the long-record sites, mask design, observed donors and current-vintage forcing. The 42 stations represent 18 connected groups, with 24 stations in Bavaria. Group resampling and leave-region-out comparisons are descriptive and do not establish independent climates. When Bavaria is excluded, the forward controlled-gap D2-minus-stochastic MSE difference changes from \SRBayernCrossOld{} in the original analysis to \SRBayernCrossNew{} under the strict-source sensitivity, in $({}^{\circ}\mathrm{C}/\mathrm{decade})^2$, on \SRBayernCrossN{} common station--repetition cases; the sign of this small difference is sensitive to the source model. The finite-record slope is distinct from an underlying climate-process trend; neither point-error reduction nor a wide total regression interval proves that missing-data uncertainty is calibrated.

On the common forward controlled-gap cohort (\UCCommon{} of \UCAll{} station--repetition cases, \UCStations{} stations in \UCGroups{} connected groups), the incomplete-year-weighted D2 90\% annual coverage is \UCAnnual\%; equal weighting of station--repetition cases gives \UCAnnualEqual\%. With matched random ranks in the tail experiment, constant versus normal-spacing tails yield \UCTailConstant\% versus \UCTailNormal\% coverage, widths \UCTailWidthConstant{} versus \UCTailWidthNormal{} and interval scores \UCTailScoreConstant{} versus \UCTailScoreNormal{} for the nominal 95\% slope interval. In the separate fixed-marginal coupling experiment, independent, auxiliary shared-parameter and residual-block ranks yield coverage \UCIndependent\%, \UCShared\% and \UCBlocks\%, with interval scores \UCIndependentScore{}, \UCSharedScore{} and \UCBlocksScore{}, respectively. The two experiments use different random draws; annual calibration has not been experimentally repaired. These contrasts are not additive causal contributions, and their cohort excludes cases where a comparison is unavailable.

The bounded annual representations also contain many degenerate 90\% intervals. On the \SPCalendarCommon{} common forward empirical-calendar cases with a finite bound, D2 has \SPCalendarDTwoZero{} zero-width intervals (\SPCalendarDTwoPercent\%) and calibrated stochastic reconstruction has \SPCalendarClassicZero{} (\SPCalendarClassicPercent\%); all of these coincide with the complete reported reference within numerical tolerance. In the \SPCrossCommon{} common forward controlled-gap cases, calibrated stochastic reconstruction has \SPCrossClassicZero{} zero-width intervals, \SPCrossClassicMiss{} of which miss the reference across \SPCrossClassicMissStations{} stations; the largest miss is \SPCrossClassicMaxMiss{}$^{\circ}$C. Among the \SPCrossDTwoMedian{} D2 cases whose median equals the bound, the reference is higher in \SPCrossDTwoMedianAbove{}; among the \SPCrossClassicMedian{} corresponding stochastic cases, it is higher in \SPCrossClassicMedianAbove{}. These are repeated masked cases, not independent rivers. Matching an observed bound can accurately recover the reported complete-record statistic when its maximum window is already visible, but does not establish zero sensor uncertainty. Replaying the fixed models and projection-aware offsets shows that, in the forward empirical-calendar cases, \SPCalendarDTwoCreatedLast{} D2 intervals first become zero-width at the final projection and \SPCalendarDTwoBeforeLast{} are already zero-width immediately before it. For the stochastic comparator, \SPCalendarClassicCreatedLast{} of its zero-width intervals first appear at this last step; \SPCalendarClassicRawZero{} raw stochastic intervals are already zero-width before calibration and the explicit projections. These stochastic quantiles come from a finite set of imputations that retains the observed daily values; raw therefore does not mean unconstrained by observations. The pre-final values still reflect the first projection and offsets calibrated with that projection; this replay does not measure the effect of refitting under an alternative calibration protocol.

An outcome-defined diagnostic separates cases in which a fully observed window reaches the complete reported H30, within $10^{-8}$ $^{\circ}$C, from cases in which the reference exceeds the observed bound. In the \PVCalendarAll{} common forward empirical-calendar cases, the bound matches the reference in \PVCalendarVisible{}; on the remaining \PVCalendarHidden{} cases, D2 and stochastic 90\% interval coverage are \PVCalendarDTwoCoverage\% and \PVCalendarClassicCoverage\%, respectively. In the \PVCrossHidden{} of \PVCrossAll{} common forward controlled-gap cases whose reference exceeds the bound, the corresponding coverages are \PVCrossDTwoCoverage\% and \PVCrossClassicCoverage\%; descriptive connected-group bootstrap intervals are [\PVCrossDTwoCILow{}, \PVCrossDTwoCIHigh{}]\% and [\PVCrossClassicCILow{}, \PVCrossClassicCIHigh{}]\%, respectively. This stratification uses the hidden complete reference and is not a decision-time diagnostic. Repeated masks can place the same station-year in both strata, and group resampling does not establish transfer across independent climates.

The candidate-provider strict-region refits were also compared on the same four-way common annual cases. Among \OpenAnnualCases{} forward controlled-gap cases at \OpenAnnualStations{} stations in \OpenAnnualGroups{} connected groups where the complete reported H30 exceeds the visible bound, D2 nominal 90\% coverage changes from \OpenAnnualDTwoOld\% to \OpenAnnualDTwoNew\%, while its mean width changes from \OpenAnnualDTwoWidthOld{} to \OpenAnnualDTwoWidthNew{}$^{\circ}$C. Calibrated stochastic coverage changes from \OpenAnnualClassicOld\% to \OpenAnnualClassicNew\%, with mean width \OpenAnnualClassicWidthOld{} to \OpenAnnualClassicWidthNew{}$^{\circ}$C. This outcome-defined stratum cannot guide a real-time decision. For the classical method, the candidate-minus-original coverage change is \OpenAnnualClassicDelta{} percentage points overall, but \OpenAnnualClassicDropFoen{} after omitting FOEN and \OpenAnnualClassicDropBayern{} after omitting Bavaria from the fixed evaluation cases. These deletions do not refit models; source composition, role assignment and strict region exclusion change together in the post-hoc refit, so no single change can be assigned as the cause. Neither version establishes reliable coverage on independent climates.

We explored an additive annual interval expansion estimated from source rule-calibration groups excluded from model fitting, the original interval calibration and the held region. It uses the equal-group-weighted empirical 90th percentile of nonnegative interval misses; a region is unsupported when fewer than three source groups contribute or the empirical correction is effectively zero. Forward source cases end before 2010. In the controlled-gap forward comparison, both methods receive a supported correction for only \AEFwdSupported{} of \AEFwdCommon{} common cases (\AEFwdPct\%), all from \AEFwdStations{} station. On those same cases, D2 coverage is \AEFwdDTwoBefore\% both before and after expansion, while mean width increases from \AEFwdDTwoWidthBefore{} to \AEFwdDTwoWidthAfter{}$^{\circ}$C. The stochastic interval likewise covers \AEFwdClassicBefore\% before and \AEFwdClassicAfter\% after, with width increasing from \AEFwdClassicWidthBefore{} to \AEFwdClassicWidthAfter{}$^{\circ}$C. Neither improvement in general coverage nor a finite-sample group guarantee follows: the empirical correction is exploratory, the source groups are few, and the supported cases are highly selected. A 90\% unrandomized rank based on one maximum nonconformity score per group requires at least nine calibration groups for a finite correction; at most \AEMaxSourceGroups{} groups are available here. Such a rank would also require appropriate group exchangeability. This experiment does not recalibrate a joint trend distribution.

A subsequent post-hoc experiment uses the candidate-source models and single-year masked cases from their rule-calibration groups to estimate an equal-group-weighted empirical 90th percentile of annual interval misses. These groups are disjoint from estimator fitting and original interval calibration, although some contributed to the separate risk-rule component; they are not a newly unopened dataset. Forward source cases precede 2010. Corrections are withheld for fewer than three contributing groups or a zero empirical radius. On \SYAllSupported{} of \SYAllEligible{} common forward controlled-gap annual cases (\SYAllStations{} stations, \SYAllGroups{} connected groups), both methods have a supported correction. On this same selected subset, D2 90\% coverage changes from \SYDTwoBase\% to \SYDTwoExpanded\%, with width \SYDTwoWidthBase{} to \SYDTwoWidthExpanded{}$^{\circ}$C and interval score \SYDTwoScoreBase{} to \SYDTwoScoreExpanded{}; the classical method changes from \SYClassicBase\% to \SYClassicExpanded\%, with width \SYClassicWidthBase{} to \SYClassicWidthExpanded{}$^{\circ}$C and interval score \SYClassicScoreBase{} to \SYClassicScoreExpanded{}. In the outcome-defined hidden-peak subset of \SYHiddenSupported{} supported cases, coverage is \SYHiddenDTwoBase\% to \SYHiddenDTwoExpanded\% for D2 and \SYHiddenClassicBase\% to \SYHiddenClassicExpanded\% for the classical method. Single-year source masks do not reproduce jointly degraded target archives. The forward rule role has at most eight groups, so a finite 90\% one-maximum-per-group rank is unavailable; the empirical correction has no finite-sample or independent-climate coverage guarantee and does not by itself identify a calibrated joint trend distribution.

A further post-hoc propagation uses the same masks and rank matrices for the source-supported annual corrections. The original forward controlled-gap sample has \SYTrendAll{} station--repetition settings; \SYTrendSupported{} have source support for both methods, and \SYTrendCommon{} nontrivial settings at \SYTrendStations{} stations in \SYTrendGroups{} connected groups have both annual distributions and all three rank constructions available. Among the supported settings, 73 lack at least one annual distribution and two contain no incomplete year. Under the independent-rank construction, D2 nominal 95\% trend coverage is \SYTrendDTwoBase\% before adjustment, \SYTrendDTwoEnd\% when only the annual 5th and 95th quantiles are shifted, and \SYTrendDTwoTaper\% when interior quantiles are also tapered outward; corresponding widths are \SYTrendDTwoWidthBase{}, \SYTrendDTwoWidthEnd{} and \SYTrendDTwoWidthTaper{}$^{\circ}$C per decade. The classical method has coverage \SYTrendClassicBase\%, \SYTrendClassicEnd\% and \SYTrendClassicTaper\%, with widths \SYTrendClassicWidthBase{}, \SYTrendClassicWidthEnd{} and \SYTrendClassicWidthTaper{}$^{\circ}$C per decade. Both adjustment schemes have identical annual 90\% endpoints, preserve the median and use the same ranks, but assume different interior and normal-spacing outer-tail shapes. The source calibration is based on single-year masks and few groups, and these retrospective results do not establish a calibrated joint distribution or independent-climate coverage.

Annual and trend coverage also depend on the assumed marginal tails. Extending the stored quantile function changes coverage even with an unchanged rank ordering, while auxiliary parameter and residual-block couplings test only particular dependence structures. A credible future joint model must account for its actual fitted-model uncertainty and retain support constraints. Task-level calibration can be useful only where independent calibration groups and relevant source difficulty are available; the small supported forward subset here cannot justify a general 95\% claim.

The decision results caution against interpreting a low average selected error as protection for every river. The imposed exposure stress is more easily detectable than unknown field faults, and measurement validity is a prerequisite for treating $L_y$ as a physical bound. Changes in groundwater influence, flow, regulation and sensor deployment can alter the water--air relationship. Their effects cannot be inferred from a model residual alone. New high-quality reference measurements are needed to evaluate natural outages and transfer beyond these archives.

\section{Conclusions}
Incomplete river-temperature records can preserve observed sustained-heat information, and retaining years reduced finite-record trend RMSE relative to the evaluated warm-season screen in the tested degradation settings. The benefit of rejecting or retaining a year depends on the actual completeness rule and the intended scientific quantity. A well-calibrated classical reconstruction remains a strong annual comparator. Forward trend intervals, three-output release rules and external transfer require their own evidence; the present results do not certify universal numerical recovery or group safety.

\section*{Open Research}
The source observations, processing code, case-level results and figure scripts require a permanent repository and citations to be supplied by the authors. A local review candidate includes processed daily values from 95 of the 96 previously eligible source stations, permitting an independent arithmetic audit of all 35 observation witnesses and their screening-rule decisions. The LUBW station is omitted pending permission, so the full eligible population is not reconstructed. The candidate also includes ten additional USGS daily API snapshots and their observation-only replay; neither check reconstructs the main models or verifies subdaily sampling. A local workflow, source inventory and unresolved data-rights details are maintained in the accompanying repository. The source includes observations from LUBW; its reuse and redistribution terms require provider-specific resolution rather than an assumed open licence. Although all whole-archive target stations have recorded reuse paths, \RightsEstimatorLUBW{} of \RightsModelCombinations{} source-model estimator roles and \RightsCalibrationLUBW{} interval-calibration roles include positive-weight LUBW cases; the access and sharing conditions of the resulting models and case-level derivatives remain unresolved. The BC published daily derivative and NEON qualification use the CC-BY author archives \cite{moore2024bc} and \cite{marzolf2025neon}; the USGS intake uses the public daily-values collection \cite{usgsogc2026}. The Nova Scotia reported-clock subset is derived from the provincial open-data release \cite{novascotiawater2026} and uses NASA POWER forcing \cite{nasapower2026}; the separately accessed legacy workbooks require their own reuse clarification. Derived files record candidate rosters and unavailable cases. Author, affiliation, funding, conflicts of interest, permitted reviewer access and permanent deposit information remain to be supplied before submission.

\input{simple_core_figures.tex}
\bibliography{references_tr,references_review_completion}
\end{document}
